% Generated draft: requires mathematical and source review.
\documentclass[11pt]{article}
\usepackage[margin=25mm]{geometry}
\usepackage{fontspec}
\setmainfont{Latin Modern Roman}
\usepackage{amsmath,amssymb,mathtools,unicode-math}
\setmathfont{Latin Modern Math}
\usepackage{enumitem,needspace,xurl,hyperref}
\setlength{\parindent}{0pt}
\setlength{\parskip}{5pt}
\setlength{\emergencystretch}{3em}
\newcommand{\N}{\mathbb{N}}
\newcommand{\Z}{\mathbb{Z}}
\newcommand{\Q}{\mathbb{Q}}
\newcommand{\R}{\mathbb{R}}
\newcommand{\C}{\mathbb{C}}
\newcommand{\F}{\mathbb{F}}
\newcommand{\A}{\mathbb{A}}
\newcommand{\PP}{\mathbb{P}}
\newcommand{\E}{\mathbb{E}}
\newcommand{\eps}{\varepsilon}
\newcommand{\dd}{\,\mathrm{d}}
\begin{document}
\section*{UnsolvedMath ID: 800310}
\textbf{Source problem code:} catalogue entry
\subsection{Definitions and mathematical statement}
Let \(S^1=\R/\Z\), and let \(f:S^1\to S^1\) be a \(C^2\) circle diffeomorphism with irrational Poincar\'e rotation number \(\rho\in\R/\Z\). Here the usual rotation-number convention implicitly requires \(f\) to preserve orientation. For a lift \(F:\R\to\R\) satisfying \(F(x+1)=F(x)+1\), this number is
\[
\rho=\left[\lim_{n\to\infty}\frac{F^n(x)-x}{n}\right]\in\R/\Z,
\]
independently of \(x\) and of the choice of lift after passage to \(\R/\Z\). Irrationality means that \(\rho\) has no rational representative.

For each \(\varepsilon>0\), define the mapping torus
\[
M_{\varepsilon}=\bigl(S^1\times[0,\varepsilon]\bigr)/\!\sim,
\qquad (x,0)\sim(f(x),\varepsilon).
\]
The supplied source then proposes a parameter \(\rho(\varepsilon)\in\R/\Z\) through the asserted correspondence
\[
M_{\varepsilon}\ \text{corresponds to}\ 
\C/\bigl(\Z+i\rho(\varepsilon)\Z\bigr),
\]
and asks whether, for every such \(f\),
\[
\lim_{\varepsilon\to0^+}\rho(\varepsilon)=\rho
\quad\text{in }\R/\Z.
\]

\textbf{Unresolved formulation qualifications.} The displayed elliptic-curve expression is retained exactly as supplied, but it does not define \(\rho(\varepsilon)\) as written. Multiplication of an element of \(\R/\Z\) by \(i\) does not specify a complex lattice generator. Choosing a real representative \(r\) would give \(\Z+ir\Z\), which depends on that choice and is not a rank-two lattice when \(r=0\). The mapping-torus construction specifies a topological surface, but the supplied record does not explain the conformal structure, the meaning of ``corresponds,'' or the marking and normalization used to extract a parameter. In the usual presentation \(\C/(\Z+\tau\Z)\), one takes a complex modulus \(\tau\) with \(\operatorname{Im}\tau>0\), with additional choices needed to select a representative. No replacement of the source's parameter or lattice formula is assumed here. Consequently, the proposed convergence question requires clarification of these data before it is a fully determined mathematical statement.
\subsection{Short English statement}
The source asks whether an elliptic-curve parameter associated with increasingly thin mapping tori of a \(C^2\) circle diffeomorphism converges to its irrational rotation number. Its stated parameter and lattice formula do not give a well-defined construction, so the intended parameter, conformal structure, and normalization require clarification. This is a statement-only draft.
\subsection{Sources}
\begin{itemize}
\item Local UnsolvedMath registry, permanent ID \(800310\), source problem code \texttt{catalogue entry}, title ``Group Actions on 1-Manifolds --- Question 27.'' The supplied record attributes the problem to Navas (2017) and labels it \texttt{open}; that label is supplied registry metadata, not an independently verified assessment.
\item Navas, \emph{Group Actions on 1-Manifolds: A List of Very Concrete Open Questions} (2017), source item recorded as ``Question 27 (\#26).'' Supplied source URL: \url{https://arxiv.org/abs/1712.06462}. The background records access on 2026-10-16 and extraction by \texttt{arxiv\_latex}; these are supplied provenance claims, not a claim of independent access or verification in this draft.
\end{itemize}
\end{document}
